\documentclass[11pt]{article}

\usepackage[margin=1in]{geometry}
\usepackage{amsmath,amssymb,amsthm}
\usepackage{hyperref}
\usepackage{xurl}
\usepackage{microtype}
\usepackage{enumitem}
\usepackage{booktabs}
\usepackage{listings}
\usepackage{float}
\usepackage[nameinlink,noabbrev]{cleveref}

\theoremstyle{definition}
\newtheorem{definition}{Definition}
\theoremstyle{remark}
\newtheorem{remark}{Remark}

\title{Closed-View Auditing and CPPCs (Addendum):\\Multi-Stream Monitoring Notes and Schema Excerpts}
\author{}
\date{March 19, 2026 (archive v393)}

\begin{document}
\maketitle

\begin{abstract}
This companion note collects only the supplementary control-plane packet for the main closed-view auditing paper (State~3):
(i) a compact many-detectors monitoring note using e-values and scan-inflation control,
(ii) a short CPPC schema excerpt,
(iii) a reference-linter excerpt, and
(iv) two illustrative CPPC card fragments.
It is intentionally narrow: we do not restate the closed-view monitoring theorem/root claim, evaluator notarization surface, runtime-plan conformance receipts, or public release-lineage packaging.
The note exists so State~3 can stay referee-tight while still pointing to one small, auditable appendix-level packet.
\end{abstract}

\paragraph{Scope.}
Addendum to State~3.\cite{series:state3}

\paragraph{Imports.}
We import the CPPC concept, the measurement-firewall pattern, and the alert-oracle framing from State~3.\cite{series:state3}
The surrounding evaluation / publication stack is imported rather than re-owned here: evaluator notarization belongs to Eval~1, runtime-plan conformance receipts belong to Operational~B, and public claim packaging / release lineage belong to ABOM/OINL plus Release~A.\cite{series:evalnotary,series:opB,series:abomoinl,series:releaseproperty}
Artifact availability follows the artifact policy.\cite{series:artifactpolicy}

\paragraph{Exports.}
A compact supplementary packet: one many-detectors monitoring note, one short CPPC schema excerpt, one reference-linter excerpt, and two illustrative example-card fragments.

\paragraph{Boundary.}
This addendum owns only those supplementary monitoring and schema artifacts.
It does \emph{not} own the closed-view monitoring contract and alert-oracle tax (State~3), evaluator notarization / anti-grinding evidence release (Eval~1), runtime-plan pinning / PTL / epoch receipts (Operational~B), or public claim packaging / release lineage (ABOM/OINL plus Release~A).\cite{series:state3,series:evalnotary,series:opB,series:abomoinl,series:releaseproperty}

\paragraph{Exported object.}
The exported object is a small support packet: one many-detector scan-inflation recipe that a referee can quote, one schema/linter excerpt that makes ``machine-checkable CPPC'' concrete, and two short illustrative cards showing how the manifest vocabulary lands on real systems without reprinting the full schema/tooling bundle.

\paragraph{Earliest sufficient stop.}
This addendum is complete once a reader can cite the many-detector recipe or the CPPC schema/card excerpt as explanatory support for a State~3 claim.
If a deployment needs the underlying closed-view theorem packet, a notarized evaluation log, a runtime conformance receipt, or a public release-lineage bundle, the reader should leave this addendum and use State~3, Eval~1, Operational~B, and ABOM/OINL plus Release~A directly.\cite{series:state3,series:evalnotary,series:opB,series:abomoinl,series:releaseproperty}

\paragraph{Later terse-paper citation posture.}
For later terse papers under space pressure, default to State~3 for the closed-view monitoring contract and public CPPC claim; cite this addendum only when the exact support packet exported here --- the many-detector monitoring note, the schema excerpt, the linter excerpt, or the illustrative card fragments --- is itself live.
If the live question is only the public monitoring theorem or alert-oracle claim, stop at State~3.
If the live question is instead ``where is the concrete schema or many-detector recipe that supports that claim?'', widen to this addendum.
If the question becomes evaluator notarization, runtime conformance, or public release-lineage comparison, leave this addendum and reopen the neighboring owners.
The maintained worked example and paired cutover rule are the general model for that stop / widen / reopen discipline.\cite{series:workedexample,series:pairedterminalcutoverrules}

\paragraph{Neighboring-layer stop map.}
Table~\ref{tab:cppc-addendum-stopmap} states the smallest handoff that keeps this addendum narrow.
A referee can stay here for the many-detector recipe, schema excerpt, linter excerpt, and illustrative card fragments, but should widen immediately once the question becomes one of theorem justification, notarized evidence, runtime-plan conformance, or public release packaging.

\begin{table}[H]
\centering
\begin{tabular}{@{}p{0.34\linewidth}p{0.58\linewidth}@{}}
\toprule
Question class & Leave this addendum for \\
\midrule
What closed-view monitoring contract or alert-oracle tax justifies the claimed witness discipline? & State~3, which owns the closed-view monitoring theorem/root claim and the alert-oracle pricing rule. \\
Which evaluator log fixes the detector set, evidence window, and anti-grinding attempt record? & Eval~1, which owns notarized evaluator manifests, gap-free attempt prefixes, and verifier reports. \\
Did the monitored deployment actually follow the pinned runtime plan and publication cadence? & Operational~B, which owns plan pinning, PTL inclusion, and epoch-conformance receipts. \\
Which public card / label / release receipt makes the monitoring claim publishable across versions? & ABOM/OINL plus Release~A, which own public claim packaging and release-lineage comparison. \\
\bottomrule
\end{tabular}
\caption{Neighboring-layer stop map for the State~3 CPPC addendum. This note owns only the supplementary control-plane packet.}
\label{tab:cppc-addendum-stopmap}
\end{table}

\paragraph{A minimal public CPPC support-packet card.}
When a later short paper only needs to answer \emph{which supplementary control-plane packet made this State~3 monitoring claim concrete?}, it should stop at the small card below plus the tiny support-packet lookup drill that follows.

\begin{table}[H]
\centering
\begin{tabular}{@{}p{0.30\linewidth}p{0.60\linewidth}@{}}
\toprule
Support-packet field & Minimal public content a later terse paper may quote here \\
\midrule
Packet claim & This addendum exports only the State~3 support packet: one many-detector monitoring note, one CPPC schema excerpt, one linter excerpt, and two illustrative card fragments. \\
Many-detector anchor & The packet's monitoring note is the e-process / e-BH scan-inflation recipe in \Cref{app:multi}; it explains how multiple closed-view detectors are monitored without turning this addendum into a second theorem paper. \\
Schema excerpt anchor & The public schema excerpt in \Cref{app:cppc} exposes the CPPC fields a later paper may name directly: \texttt{witness\_channels\_present}, \texttt{cadence}, \texttt{view\_accountant}, \texttt{budget\_beacon}, \texttt{transparency}, and \texttt{compliance}. \\
Linter anchor & The \texttt{tools/cppc\_lint.py} excerpt shows the machine-checkable support rule: required keys present, enum check on witness channels, and artifact-existence / hash checks when provided. \\
Illustrative card anchors & The two public card fragments are the Nym-style and Tor-style excerpts below; they are explanatory support for the CPPC vocabulary, not new theorem or release objects. \\
How to cite tersely & Quote this card only when the live issue is the concrete support packet itself; otherwise stop at State~3 for the public monitoring claim. \\
Do not widen silently & Evaluator notarization, runtime-plan conformance, and public release lineage still belong to Eval~1, Operational~B, and ABOM/OINL plus Release~A. \\
\bottomrule
\end{tabular}
\caption{A minimal public CPPC support-packet card. This is the smallest public answer later terse papers can quote directly when the live issue is the concrete many-detector note, schema excerpt, linter excerpt, or illustrative CPPC card fragments rather than the main State~3 monitoring contract.}
\label{tab:cppc-addendum-card}
\end{table}

\paragraph{Tiny support-packet lookup drill.}
Suppose a later short paper only needs to justify the sentence ``the monitored deployment published a CPPC support packet with \texttt{budget\_beacon} and \texttt{auditing\_alerts} channels, an RDP odometer, and a machine-checkable linter path.''
The smallest public support object is exactly this addendum card plus the schema excerpt in \Cref{app:cppc} and the Nym-style illustrative fragment below.
That packet already names \texttt{witness\_channels\_present = [budget\_beacon, auditing\_alerts, timing\_and\_metadata]}, a one-minute epoch / daily delayed publication cadence, and the reference linter checks.
A later terse paper can therefore cite one concrete CPPC support packet without reopening State~3's theorem root, Eval~1's notarized evidence packet, Operational~B's runtime receipts, or Release~A's lineage object.

\appendix


\section{Multi-stream monitoring note (e-values)}
\label{app:multi}
Stateful anonymity deployments often monitor many candidate witnesses (many committees/keys/relays/metrics).
A concise way to control ``scan inflation'' is to attach each detector an \emph{e-process} (a nonnegative supermartingale with unit expectation under the null) and apply e-value multiple testing.

\paragraph{E-values and e-BH.}
E-values generalize likelihood-ratio betting scores; they compose naturally under optional stopping and enable dependence-robust FDR control via the e-BH procedure~\cite{ramdas2024evalues,wang2022ebh}.
Operationally, this lets an operator run many closed-view detectors while bounding the expected false discovery proportion.

\paragraph{Stopping-time subtlety.}
When alerts are raised at data-dependent global stopping times (``whenever anything looks bad''), additional care is required; the stopped e-BH procedure is designed for this regime~\cite{wang2025sebH}.
The main paper develops the deployment-side implication: \emph{publishing} the alert time/window can itself be a witness stream and should be priced.\cite{series:state3}



\section{CPPC schema (excerpt) and example cards}
\label{app:cppc}
To keep this note referee-tight, we include only a short \emph{excerpt} of the CPPC schema and two illustrative cards.
The full JSON schema, complete card examples, and machine-checking tooling remain available on request as part of the series archive under the repository artifact policy.\cite{series:artifactpolicy}

\subsection{Schema excerpt (v1.1)}
\begin{lstlisting}
{
  "schema_version": "1.1",
  "system": "string",
  "unit_of_privacy": { "...": "..." },
  "witness_streams": [ "string", "..." ],
  "witness_channels_present": [
    "semantic_outputs",
    "budget_beacon",
    "transparency_log",
    "timing_and_metadata",
    "auditing_alerts"
  ],
  "cadence": { "epoch": "...", "publication": "..." },
  "view_accountant": {
    "notion": "max-divergence | RDP | f-DP/GDP",
    "composition": "...",
    "budget": { "...": "..." }
  },
  "budget_beacon": { "cadence": "...", "format": "..." },
  "transparency": { "log": "...", "signatures": "..." },
  "compliance": { "mode": "...", "artifacts": [ "..." ] }
}
\end{lstlisting}

\subsection{Reference linter (excerpt)}
In the draft archive we include a small reference linter (\texttt{tools/cppc\_lint.py}) that checks a CPPC JSON for required fields and basic internal consistency.
The intent is to make ``machine-checkable'' concrete without requiring any particular SBOM or transparency-log ecosystem.

\begin{lstlisting}
$ python tools/cppc_lint.py card.json --root ./releases/
OK: required keys present
OK: witness_channels_present is a subset of allowed enum
OK: compliance.artifacts exist under --root and match declared sha256 (if provided)
WARN: transparency.log missing (deployment-specific)
WARN: transparency.signatures missing (deployment-specific)
\end{lstlisting}

In the stop map above, that linter remains only a support artifact: it makes the CPPC vocabulary auditable, but it does not by itself settle whether the monitored deployment satisfied the State~3 theorem assumptions, matched the pinned runtime plan, or was packaged into a publishable release claim.



\subsection{Example CPPC (Nym-style, illustrative excerpt)}
\begin{lstlisting}
{
  "schema_version": "1.1",
  "system": "nym-mixnet (illustrative)",
  "unit_of_privacy": { "message": "1 packet-equivalent" },
  "witness_channels_present": [
    "budget_beacon",
    "auditing_alerts",
    "timing_and_metadata"
  ],
  "cadence": {
    "epoch": "1 min",
    "publication": "daily aggregates; delayed by 7d"
  },
  "witness_streams": [
    "mix-queue depth (bucketed)",
    "node reputation updates (coarsened)",
    "drift alerts (delayed, bucketed)"
  ],
  "view_accountant": {
    "notion": "RDP",
    "composition": "odometer (adaptive cadence)",
    "budget": { "epsilon_rdp": 3.0, "delta": 1e-6, "horizon": "30d rolling" }
  },
  "budget_beacon": {
    "cadence": "daily",
    "payload": "current spend + remaining budget (coarsened)"
  },
  "transparency": {
    "log": "append-only merkle log",
    "signatures": "operator key + optional attestation"
  },
  "compliance": {
    "mode": "attestation",
    "artifacts": [ "cppc.json", "build hash", "release hashes" ]
  }
}
\end{lstlisting}



\subsection{Example CPPC (Tor-style, illustrative excerpt)}
\begin{lstlisting}
{
  "schema_version": "1.1",
  "system": "tor-relay-selection (illustrative)",
  "unit_of_privacy": { "client": "one long-lived guard set" },
  "witness_channels_present": [
    "semantic_outputs",
    "budget_beacon",
    "transparency_log"
  ],
  "cadence": {
    "epoch": "1 consensus period",
    "publication": "public consensus + coarse metrics"
  },
  "witness_streams": [
    "selection weights (public consensus)",
    "bandwidth measurement summaries (coarse)",
    "guard selection randomness"
  ],
  "view_accountant": {
    "notion": "max-divergence",
    "composition": "filter (abort-safe)",
    "budget": { "epsilon": 2.0, "delta": 1e-6, "horizon": "90d" }
  },
  "budget_beacon": {
    "cadence": "each consensus period",
    "payload": "monotone spend counter + checkpoint hash"
  },
  "transparency": {
    "log": "public merkle log (checkpointed)",
    "signatures": "authority keys"
  },
  "compliance": {
    "mode": "hybrid",
    "artifacts": [ "cppc.json", "consensus hashes", "auditor scripts (on request)" ]
  }
}
\end{lstlisting}







\begin{thebibliography}{99}

\bibitem{series:state3}
Anonymous.
\newblock Closed-View Auditing for Stateful Anonymity: Measurement Firewalls and CPPCs.
\newblock AnonDHT State series Paper~3, 2026. (This archive.)

\bibitem{series:contractobjects}
Anonymous.
\newblock Contract Objects for Anonymous-DHT Privacy Budgets and Claims.
\newblock Series synthesis note (Synthesis~0), 2026. (This archive.)

\bibitem{series:abomoinl}
Anonymous.
\newblock ABOM/OINL: Audit BOM and Outcome Identification Noise Label for Anonymous DHT Deployments.
\newblock Anonymity series Paper~C, 2026. (This archive.)

\bibitem{series:evalnotary}
Anonymous.
\newblock Notarized Anonymity Certificates Under Retries.
\newblock Evaluation series Paper~1, 2026. (This archive.)

\bibitem{series:opB}
Anonymous.
\newblock Auditable Trace Privacy for Anonymous DHTs: Runtime Conformance, Plan Pinning, and Epoch Receipts.
\newblock Operational series Paper~B, 2026. (This archive.)

\bibitem{series:releaseproperty}
Anonymous.
\newblock Release Property Ledger Receipts for Anonymity Claims.
\newblock Release \& destination Paper~A, 2026. (This archive.)

\bibitem{series:workedexample}
Anonymous.
\newblock Worked Example: Instantiating a Tiered Reader-Privacy Receipt.
\newblock Synthesis~17, The-One-True-Archive (2026).

\bibitem{series:pairedterminalcutoverrules}
Anonymous.
\newblock Paired Terminal Cutover Rules for Stable Review Spines: When Later Terse Papers Must Stop, Widen, or Reopen.
\newblock Series synthesis note (Synthesis~75), 2026. (This archive.)

\bibitem{series:artifactpolicy}
Anonymous.
\newblock Artifact policy and supporting materials for the anonymity/AnonDHT series.
\newblock 2026.

\bibitem{ramdas2024evalues}
Aaditya Ramdas and Ruodu Wang.
\newblock Hypothesis testing with e-values.
\newblock arXiv:2410.23614, 2024. DOI: 10.48550/arXiv.2410.23614.
\newblock \url{https://arxiv.org/abs/2410.23614}

\bibitem{wang2022ebh}
Ruodu Wang and Aaditya Ramdas.
\newblock False discovery rate control with e-values.
\newblock \emph{Journal of the Royal Statistical Society: Series B}, 84(3), 822--852, 2022.
\newblock DOI: 10.1111/rssb.12489.

\bibitem{wang2025sebH}
Hongjian Wang, Sanjit Dandapanthula, and Aaditya Ramdas.
\newblock Anytime-valid FDR control with the stopped e-BH procedure.
\newblock arXiv:2502.08539, 2025. DOI: 10.48550/arXiv.2502.08539.
\newblock \url{https://arxiv.org/abs/2502.08539}
\end{thebibliography}

\end{document}